\begin{exercise}{31}
Suppose that $X$ is a composer with the inner $3$-horn extension property. Let
$\sigma\colon\Delta^1\times\Delta^1\to X$ be a map such that
\begin{enumerate}[label=(\arabic*),leftmargin=*,itemsep=0pt,parsep=0pt,
                  topsep=.5\baselineskip]
 \item $\sigma_{\Delta^1\times\{0\}}=f$,
 \item $\sigma_{\Delta^1\times\{1\}}=g$,
 \item $\sigma_{\{0\}\times\Delta^1}=\mathrm{id}_x$, and
 \item $\sigma_{\{1\}\times\Delta^1}=\mathrm{id}_y$,
\end{enumerate}
for morphisms $f,g\colon x\to y$. Show that $f\sim g$ in the sense of
Definition~1.2.3.

\begin{proof}
Let $e\colon x\to y$ be the image of the diagonal.
The triangles with vertex lists
$((0,0),(1,0),(1,1))$ and $((0,0),(0,1),(1,1))$
map to $2$-simplices $A,B$ of $X$ with
\[
 \begin{aligned}
  (d_2A,d_1A,d_0A)&=(f,e,\mathrm{id}_y),\\
  (d_2B,d_1B,d_0B)&=(\mathrm{id}_x,e,g).
 \end{aligned}
\]
The first witnesses $f\sim e$ by Definition~1.2.3.
The hypotheses on $X$ are exactly those of Lemma~1.2.9;
part~(4), applied to $B$, gives $e\sim g$.
Part~(2) gives transitivity, hence $f\sim g$.
\end{proof}
\end{exercise}
